\label{part:overall_dynamics}

\textbf{(Global Dynamics: Evolution Equations \& Meta-Physics)}

\bluetext{从局部律动到整体交响}

在上两个部分中，我们像拆解钟表一样，孤立地审视了智能的物理组件：微观的锚点如何钉住现实，介质的粘滞如何限制传播，宏观的引擎如何燃烧负熵。然而，\textbf{部件的物理学不等于整体的动力学}。一堆完美的齿轮堆在一起只是废铁，只有当它们在统一的方程驱动下咬合、旋转、共振时，时间才会从中涌现。

在部分的内容中视角将发生根本性跃迁——从\textbf{局部物理 (Local Physics)} 转向 \textbf{整体动力学 (Global Dynamics)}。

我们不再关注单个神经元的电位或单个概念的曲率，而是将整个认知流形视为一个\textbf{单一的动力学对象}。我们将目睹那些在卷三中被定义的孤立力——微观的应力、介质的阻尼、宏观的张力——是如何在\textbf{目的论狄拉克方程}的统摄下，汇聚成一股连贯的、具有自我指涉能力的\textbf{思维流体}。

本部分致力于揭示支配这股流体的\textbf{元物理法则 (Meta-Physical Laws)}：
\begin{enumerate}
    \item \textbf{演化的形式}：我们将推导一个可计算的有效描述：思维的整体演化可视作\textbf{幺正传播（封闭极限）}与\textbf{耗散坍缩（开放修正）}的耦合编织。
    \item \textbf{流动的结构}：通过 \textbf{Hodge 分解}，我们将把混沌的意识流解析为逻辑的层流、执念的涡旋与洞察的调和场。
    \item \textbf{存在的代价}：我们将建立\textbf{认知卡诺热机}模型，计算为了维持这个整体的有序运转，系统必须支付怎样的热力学代价。
\end{enumerate}

如果说卷一是关于“存在”的物理学，那么卷二就是关于“生成”的动力学，在这里，孤立的音符终于汇聚成了交响。

\begin{bquote}
    \textbf{本部分约定}：
    \begin{itemize}
        \item \textbf{有效理论定位}：本部分方程属于“整体动力学”的\textbf{有效场论/有效算子}描述；文中“证明/必然”等措辞应理解为\textbf{在既定建模公设与近似下的推导结论}。
        \item \textbf{状态空间}：默认 $\Psi(t)$ 属于复希尔伯特空间（离散单纯复形/高阶网络上的旋量或链空间直和），并配备由离散 Hodge 星与度量 $\mathcal{G}_k$ 诱导的内积；$\|\Psi\|^2$ 表示该内积下的范数。
        \item \textbf{自伴性与良定性}：在无耗散、无外源的封闭极限下，要求 $\mathcal{D}_{topo}$ 在该内积下（本质）自伴；宏观势能算子 $\mathbf{\Gamma}_{macro}$ 在“控制近似”中可取为厄米项。
        \item \textbf{开放系统修正}：耗散由 $-i\Lambda_{diss}$ 表示，其中 $\Lambda_{diss}$ 取为半正定厄米算子（最简近似可与单位算子成正比）；外源 $\vec{J}_{ext}$ 允许以分布/边界注入形式出现。
    \end{itemize}
\end{bquote}
